% Regression test: natbib's (and biblatex's) two-optional-argument citation
% form \cite[pre][post]{key} with \cite capture active. The capture used to
% peek for a single optional argument only, so the mandatory-argument scan
% grabbed the bare token "[" as the key and typeset the rest of the citation
% as text, breaking compilation with "Missing $ inserted" (from the underscore
% in the key) deep inside the proof body. The post-note must be recorded as
% the node's note (matching the one-argument semantics) and the pre-note
% ignored: \cite[see][Section 3]{k} and \cite[Section 3]{k} must share one
% slot, as must \cite[see][]{k2} and \cite{k2}. Also exercises the capture of
% \citet and \citep (which coincide with \cite on the same (key, note) pair),
% the starred \citep* (star affects rendering only, not slot identity), and a
% multi-key \citep. Uses natbib with manual author-year \bibitem entries;
% single-pass and bibtex-free (citelabel=key) for determinism.
\documentclass{article}
\usepackage{amsthm}
\usepackage{natbib}
\usepackage[cite,citelabel=key]{proofgraph}
\newtheorem{theorem}{Theorem}
\begin{document}

\begin{theorem}\label{thm:main}Main.\end{theorem}
\begin{proof}
  With pre- and post-note, \cite[see][Section 3]{Lachal_Class_2011}, then the
  same result cited post-note only, \cite[Section 3]{Lachal_Class_2011},
  reusing its slot.
  With a pre-note only, \cite[see][]{Other_Work_2020}, then the same work
  cited plainly, \cite{Other_Work_2020}, reusing its slot too.
  A textual \citet[Section 4]{Lachal_Class_2011} opens a new slot (new note),
  while \citep[cf.][Section 3]{Lachal_Class_2011} reuses the first slot and
  the starred \citep*{Other_Work_2020} the second.
  A multi-key \citep{Lachal_Class_2011,Third_Item_1999} adds one slot per key.
\end{proof}

\begin{thebibliography}{9}
\bibitem[Lachal(2011)]{Lachal_Class_2011}A reference.
\bibitem[Other(2020)]{Other_Work_2020}Another reference.
\bibitem[Third(1999)]{Third_Item_1999}A third reference.
\end{thebibliography}
\end{document}
